% ---------------------------------------------------------------------------
% ACT D -- MULTIPOINT SHAPE OPTIMISATION (INCOMPRESSIBLE) -- one-page
% customer-facing result sheet.  Written under the DEMO STANDARD v2 output
% rules R1-R10.
%
% Self-contained: no external figure file, no packages beyond a base TeX Live
% install (no tikz/pgfplots on this box).  All three figures are native LaTeX
% picture environments, so the output is vector at any zoom.
% Compile: pdflatex ACT_D_shape_optimisation_sheet.tex
%
% PROVENANCE OF EVERY USER-VISIBLE NUMBER (not rendered; for the lab only).
%   Angle sweep, coefficients, residuals:
%     /home/ubuntu/certonomous-runs/CURRICULUM-SO3aR2-a1-naca0012-alpha-multipoint-gradient
%       /X-P_20260831T225359Z_614905.log  lines 1563-1564, 1693-1694, 1823-1824
%       residual lines 1561, 1691, 1821 ("satisfied the prescribed tolerance 1e-08")
%   Reproduction (second, earlier, independent run of the same three angles):
%     /home/ubuntu/certonomous-runs/CURRICULUM-SO3aF-a1-naca0012-alpha-feasibility
%       /out/primal_alpha_{3,5,7}.13918623195176.log lines 564-567
%     Largest relative difference over all six values: 6.6129e-09 (C_L, 5.139 deg);
%     the other four values are bit-identical between the two runs.
%   Derivatives, finite-difference table, step-size sweep, assembly identity,
%   wrong-step control:
%     .../CURRICULUM-SO3aR2-.../SO3aR2_grade_20260831T230221Z.json
%       gates.G5J.PATCHED.{G5J_objective,G5C_lift_per_scenario,
%                          G_MP_STRUCT_assembly_identity,G_TB_trivial_baseline}
%   Mesh cell count: .../MESH_20260831T224345Z_603001.log ("cells: 4032")
%   Compute used per stage: .../CURRICULUM-SO3aR2-.../ledger.txt
%     MESH 0.167 + X-S 2.517 + F-S 4.367 + X-P 2.517 + F-P 4.75 = 14.318 core-min
%     wall 10 + 151 + 262 + 151 + 285 = 859 s
%   Feasibility-check wall/core-min/cap -- A DIFFERENT RUN ROOT, and it has no
%   ledger.txt; its figures come from its own status file:
%     /home/ubuntu/certonomous-runs/CURRICULUM-SO3aF-a1-naca0012-alpha-feasibility
%       /STATUS.SO3aF   -> wall_s=35 ranks=1 core_min=0.5833 cap_core_min=8.0
%   Upfront estimate 22.1 core-min and the 5 per cent acceptance limit:
%     cases/dafoam/ladder-a/A1/curriculum_SO3aR2/PREREGISTRATION.md lines 75-77, 226
%   Case setup (U0=10, nu=1.5e-5, A0=0.1, chord 1 m, 5x2x2 shape controls,
%   thickness/volume/radius constraints):
%     .../CURRICULUM-SO3aR2-.../so3ar2_runScript.py lines 143, 151, 331-344,
%     379-381; MESH/constant/transportProperties; MESH/FFD/wingFFD.xyz
%
% FIGURE SPECIFICATIONS (all three drawn below in native picture units,
% 1 unit = 1 mm; these notes state what each drawing means and how it is built).
%
%   FIG 1  DESCENT DIRECTION ON THE SECTION -- the money visual.
%     Geometry: NACA0012 to true scale, chord 1 m -> 62 mm, x to the right from
%     the leading edge, y vertical, EQUAL ASPECT (no thickness exaggeration).
%     Axes: x is chord station x/c, dimensionless; y is section thickness in the
%     same length units.  No y axis is drawn because the section is to scale.
%     Arrows: the NEGATIVE of the verified derivative of the weighted-mean drag
%     coefficient with respect to each surface control -- i.e. the direction
%     that LOWERS drag -- normalised so the largest is 7.5 mm.
%       control at x/c=0.245, lower surface : dJ/ds = -0.012248354486 m^-1
%           -> descent +y, 5.51 mm, one arrow (single-surface control)
%       control at x/c=0.500, upper surface : dJ/ds = +0.012013942626 m^-1
%           -> descent -y, 5.41 mm, one arrow (single-surface control)
%       leading-edge thickness control       : dJ/ds = +0.016660787054 m^-1
%           -> descent = thicker, 7.50 mm, drawn as a PAIR (this control moves
%              upper and lower surfaces symmetrically)
%       trailing-edge thickness control      : dJ/ds = +0.006202686046 m^-1
%           -> descent = thicker, 2.79 mm, drawn as a PAIR
%     Colour: blue = single-surface control, orange = symmetric pair.
%     Scale bar: a 7.5 mm arrow = 0.0167 per metre of control travel.
%     Only the four controls checked against finite differences are drawn; the
%     other four shape controls are not verified and are not shown.
%
%   FIG 2  STEP-SIZE CHECK.
%     x: finite-difference step, logarithmic, 1e-4 to 1e-2 (metres of control
%        travel); 1e-4 -> 0 mm, 1e-3 -> 20 mm, 1e-2 -> 40 mm.
%     y: (finite-difference estimate - computed derivative) / |computed
%        derivative| x 100, from -5 to +11 per cent; y_mm = (percent + 5)*1.625.
%     Shaded band: the +/-5 per cent acceptance limit fixed before the run.
%     Four lines, one per checked control; values (h=1e-4, 1e-3, 1e-2, per cent):
%        x/c=0.245 lower   : +0.5847, +0.0464, +0.5315
%        x/c=0.500 upper   : +0.6258, +0.0418, +0.0937
%        leading-edge thick: -0.2733, -3.9012, +10.2563
%        trailing-edge thick: +1.2476, +0.1448, +0.2215
%
%   FIG 3  POLAR.
%     x: angle of attack, 3.0 to 7.5 deg -> 0 to 38 mm.
%     Left axis: lift coefficient 0.30-0.70 -> 0 to 24 mm, filled circles,
%       solid line.  Right axis: drag coefficient 0.015-0.030 -> 0 to 24 mm,
%       open squares, dashed line.  Three solved points on each curve.
%     Error bars are smaller than the markers (Table 2) and are not drawn.
%     No experimental reference curve exists for this section at this
%     condition, so none is plotted.
% ---------------------------------------------------------------------------
\documentclass[9pt,a4paper]{article}
\usepackage[margin=8mm,top=6mm,bottom=3mm]{geometry}
\usepackage{booktabs}
\usepackage{array}
\usepackage{xcolor}
\usepackage{graphicx}
\usepackage[T1]{fontenc}

\definecolor{band}{rgb}{0.87,0.91,0.96}
\definecolor{accent}{rgb}{0.10,0.32,0.60}
\definecolor{warm}{rgb}{0.72,0.30,0.06}
\definecolor{rulegrey}{rgb}{0.45,0.45,0.45}
\definecolor{boxgrey}{rgb}{0.96,0.96,0.94}
\definecolor{foil}{rgb}{0.78,0.78,0.76}

\renewcommand{\footnotesize}{\fontsize{6.4}{7.05}\selectfont}
\renewcommand{\scriptsize}{\fontsize{5.7}{6.3}\selectfont}
\renewcommand{\arraystretch}{1.03}
\setlength{\parindent}{0pt}
\setlength{\parskip}{0.5pt}
\setlength{\tabcolsep}{3.2pt}
\pagestyle{empty}

\newcommand{\rolesig}[1]{\textbf{\footnotesize#1}}

\begin{document}
\fontsize{7.3}{8.5}\selectfont

% ---------------------------------------------------------------- title -----
{\fontsize{13.5}{15}\selectfont\bfseries Multipoint shape optimisation of an
aerofoil section}\\[1.5pt]
{\fontsize{8.6}{10}\selectfont The gradients that will drive it, verified at
three angles of attack --- incompressible, steady, two-dimensional}\\[2pt]
{\color{rulegrey}\rule{\textwidth}{0.9pt}}\\[1pt]
{\footnotesize \textbf{Shape optimisation runs next.} This sheet is the step
before it: the derivatives that tell the optimiser which way to move the
surface, checked against an independent calculation. No optimisation iteration
has been taken, so there is no shape change and no drag-reduction figure here.}

% --------------------------------------------- assumption beat + table 1 ----
\vspace{2.5pt}
\begin{minipage}[t]{0.505\textwidth}
\textbf{One thing we changed in your request}

\vspace{1pt}
{\footnotesize You asked for the shape to be optimised across three Mach
numbers. At ten metres per second the flow is incompressible --- Mach $\approx
0.03$ (Table 1), where compressibility corrections are a fraction of a per
cent. Mach is not a parameter of this problem: three Mach points would be three
copies of the same run, and you would pay three times for one answer. Angle of
attack is where the shape actually trades off, so that is the sweep to optimise
across.\\[2pt]
\textbf{The problem we set up.} Minimise the drag coefficient averaged over
three angles of attack, weighted equally, with one shared set of surface
controls serving all three --- so a single shape must work across the range
rather than being tuned to one point. Minimum thickness, minimum enclosed area
and minimum leading-edge radius are held as constraints. The angles, the
weights and the acceptance limit on the derivative check were all written down
and frozen before the solver started.}
\end{minipage}\hfill
\begin{minipage}[t]{0.475\textwidth}
\textbf{Table 1 --- Operating condition.}

\vspace{1.5pt}
{\fontsize{6.7}{7.9}\selectfont
\begin{tabular}{@{}l r l >{\raggedright\arraybackslash}p{22mm}@{}}
\toprule
Quantity & Value & Unit & Uncertainty / basis \\
\midrule
Freestream speed $U$        & $10$                 & m\,s$^{-1}$    & exact input \\
Chord $c$                   & $1$                  & m              & exact input \\
Viscosity $\nu$ (kinematic) & $1.5\times10^{-5}$   & m$^2$\,s$^{-1}$& exact input \\
Reynolds number $Uc/\nu$    & $6.67\times10^{5}$   & ---            & derived, exact \\
Mach number $U/a$           & $0.029$              & ---            & derived, $a=340$\,m\,s$^{-1}$ \\
$1/\sqrt{1-M^{2}}$          & $1.0004$             & ---            & textbook \\
Angles of attack $\alpha$   & $3.139$, $5.139$, $7.139$ & deg       & fixed beforehand \\
Mesh                        & $4\,032$             & cells          & counted \\
\bottomrule
\end{tabular}}

\vspace{1pt}
{\scriptsize $a$ is the speed of sound in standard sea-level air. The Mach
row is a property of the stated speed and of standard air, not a finding of
this work: nothing on this sheet is a compressibility measurement.}
\end{minipage}

% ------------------------------------ fig 1 + table 3  |  table 2 + fig 2 ---
\vspace{1pt}
\begin{minipage}[t]{0.545\textwidth}
\textbf{Figure 1 --- Which way the optimiser would move the surface.}
Arrows point in the direction that lowers the averaged drag; length is
proportional to how strongly. Drawn from the four verified derivatives in
Table 3, section at true scale.

\vspace{1.5pt}
\setlength{\unitlength}{1mm}
\begin{picture}(83,23)(-5,-11.5)
  % ---- chord line ---------------------------------------------------------
  \linethickness{0.4pt}
  {\color{foil}\multiput(0,0)(3,0){21}{\rule{1.6mm}{0.25pt}}}
  % ---- aerofoil upper surface (NACA0012, chord 62 mm) ---------------------
  \linethickness{0.55pt}
  \qbezier(0.00,0.00)(0.09,0.49)(0.31,0.75)
  \qbezier(0.31,0.75)(0.53,0.99)(0.78,1.18)
  \qbezier(0.78,1.18)(1.15,1.43)(1.55,1.62)
  \qbezier(1.55,1.62)(2.30,1.96)(3.10,2.21)
  \qbezier(3.10,2.21)(4.65,2.59)(6.20,2.90)
  \qbezier(6.20,2.90)(7.75,3.13)(9.30,3.31)
  \qbezier(9.30,3.31)(10.85,3.45)(12.40,3.56)
  \qbezier(12.40,3.56)(15.50,3.69)(18.60,3.72)
  \qbezier(18.60,3.72)(21.70,3.69)(24.80,3.60)
  \qbezier(24.80,3.60)(27.90,3.46)(31.00,3.29)
  \qbezier(31.00,3.29)(34.10,3.07)(37.20,2.82)
  \qbezier(37.20,2.82)(40.30,2.56)(43.40,2.27)
  \qbezier(43.40,2.27)(46.50,1.96)(49.60,1.63)
  \qbezier(49.60,1.63)(52.70,1.28)(55.80,0.89)
  \qbezier(55.80,0.89)(57.35,0.70)(58.90,0.50)
  \qbezier(58.90,0.50)(60.45,0.29)(62.00,0.08)
  % ---- aerofoil lower surface ---------------------------------------------
  \qbezier(0.00,0.00)(0.09,-0.49)(0.31,-0.75)
  \qbezier(0.31,-0.75)(0.53,-0.99)(0.78,-1.18)
  \qbezier(0.78,-1.18)(1.15,-1.43)(1.55,-1.62)
  \qbezier(1.55,-1.62)(2.30,-1.96)(3.10,-2.21)
  \qbezier(3.10,-2.21)(4.65,-2.59)(6.20,-2.90)
  \qbezier(6.20,-2.90)(7.75,-3.13)(9.30,-3.31)
  \qbezier(9.30,-3.31)(10.85,-3.45)(12.40,-3.56)
  \qbezier(12.40,-3.56)(15.50,-3.69)(18.60,-3.72)
  \qbezier(18.60,-3.72)(21.70,-3.69)(24.80,-3.60)
  \qbezier(24.80,-3.60)(27.90,-3.46)(31.00,-3.29)
  \qbezier(31.00,-3.29)(34.10,-3.07)(37.20,-2.82)
  \qbezier(37.20,-2.82)(40.30,-2.56)(43.40,-2.27)
  \qbezier(43.40,-2.27)(46.50,-1.96)(49.60,-1.63)
  \qbezier(49.60,-1.63)(52.70,-1.28)(55.80,-0.89)
  \qbezier(55.80,-0.89)(57.35,-0.70)(58.90,-0.50)
  \qbezier(58.90,-0.50)(60.45,-0.29)(62.00,-0.08)
  % ---- arrows: the direction that lowers drag ------------------------------
  \linethickness{0.6pt}
  % leading-edge thickness control -- symmetric pair, longest (7.50 mm)
  \put(1.86,2.05){{\color{warm}\vector(0,1){7.50}}}
  \put(1.86,-2.05){{\color{warm}\vector(0,-1){7.50}}}
  \put(0,11.3){\makebox(0,0)[l]{\scriptsize\color{warm}leading edge}}
  % lower surface at x/c = 0.245 -- upward (5.51 mm)
  \put(15.19,-3.68){{\color{accent}\vector(0,1){5.51}}}
  \put(13.6,-6.6){\makebox(0,0)[c]{\scriptsize\color{accent}$x/c=0.245$}}
  % upper surface at x/c = 0.500 -- downward (5.41 mm)
  \put(31.00,3.29){{\color{accent}\vector(0,-1){5.41}}}
  \put(31.00,5.3){\makebox(0,0)[c]{\scriptsize\color{accent}$x/c=0.500$}}
  % trailing-edge thickness control -- symmetric pair (2.79 mm)
  \put(60.14,0.35){{\color{warm}\vector(0,1){2.79}}}
  \put(60.14,-0.35){{\color{warm}\vector(0,-1){2.79}}}
  \put(62.5,-6.6){\makebox(0,0)[r]{\scriptsize\color{warm}trailing edge}}
  % ---- scale bar ----------------------------------------------------------
  \linethickness{0.5pt}
  \put(66,1){\line(0,1){7.5}}
  \put(65.3,1){\line(1,0){1.4}}
  \put(65.3,8.5){\line(1,0){1.4}}
  \put(67.4,4.75){\makebox(0,0)[l]{\scriptsize $0.0167\,\mathrm{m^{-1}}$}}
  % ---- axis label ---------------------------------------------------------
  \put(31,-10.2){\makebox(0,0)[c]{\scriptsize chord station $x/c$ --- section
      to true scale, chord $c=1\,\mathrm{m}$}}
\end{picture}

\vspace{0.5pt}
{\scriptsize Blue arrows are single-surface controls; orange pairs move both
surfaces, and at both edges the objective wants the section \emph{thicker}.
This is the local direction of steepest drag reduction at your current shape
--- the thickness, area and radius constraints will bend the actual step.}

\vspace{2pt}
\textbf{Table 3 --- The derivative check.} Each surface control's effect on the
averaged drag, computed two independent ways.

\vspace{1.5pt}
{\fontsize{6.7}{7.9}\selectfont
\begin{tabular}{@{}l r r r@{}}
\toprule
Surface control & Computed, m$^{-1}$ & Finite difference, m$^{-1}$ &
  \begin{tabular}{@{}r@{}}Difference,\\ \% of computed\end{tabular} \\
\midrule
Lower surface, $x/c=0.245$ & $-0.012248$ & $-0.012243$ & $0.046$ \\
Upper surface, $x/c=0.500$ & $+0.012014$ & $+0.012019$ & $0.042$ \\
Leading-edge thickness     & $+0.0167$   & $+0.0160$   & $\mathbf{4.06}$ \\
Trailing-edge thickness    & $+0.006203$ & $+0.006212$ & $0.145$ \\
\midrule
\multicolumn{3}{@{}l}{Acceptance limit, fixed before the run} & $5.00$ \\
\bottomrule
\end{tabular}}

\vspace{1pt}
{\scriptsize Each value carries one digit beyond its own difference column, so
no digit is printed that the check does not support. Units are per metre of
control travel; chord is $1\,\mathrm{m}$, so this is also per chord. Lift
derivatives were checked the same way at all three angles: worst disagreement
$1.185\,\%$, best $0.004\,\%$.}
\end{minipage}\hfill
\begin{minipage}[t]{0.435\textwidth}
\textbf{Table 2 --- The three solved operating points.}

\vspace{1.5pt}
{\fontsize{6.7}{7.9}\selectfont
\begin{tabular}{@{}r r r l@{}}
\toprule
$\alpha$, deg & $C_D$, --- & $C_L$, --- &
  \begin{tabular}{@{}l@{}}Run-to-run reproduction,\\ relative\end{tabular} \\
\midrule
$3.139$ & $0.0172394$ & $0.311896$ & $0$ (identical to last digit) \\
$5.139$ & $0.0209105$ & $0.498765$ & $6.6\times10^{-9}$ \\
$7.139$ & $0.0272681$ & $0.663976$ & $0$ (identical to last digit) \\
\bottomrule
\end{tabular}}

\vspace{1pt}
{\scriptsize These three points were solved twice, in two separate runs made
hours apart; the largest relative difference across all six values is
$6.6\times10^{-9}$. That column is repeatability, not accuracy. Drag rises
monotonically across the range, and lift at the middle angle is $0.498765$
against your $0.500$ design target.}

\vspace{2pt}
\textbf{Figure 2 --- Step-size check.} The finite-difference estimate against
the computed derivative, as the difference step is varied.

\vspace{1.5pt}
\setlength{\unitlength}{1mm}
\begin{picture}(58,30)(-8,-5.5)
  \put(0,0){{\color{band}\rule{40mm}{16.25mm}}}
  \linethickness{0.4pt}
  \put(0,0){\line(1,0){40}}
  \put(0,0){\line(0,1){27}}
  {\color{rulegrey}\multiput(0,8.125)(3,0){14}{\rule{1.7mm}{0.3pt}}}
  % y ticks: -5, 0, +5, +10
  \multiput(0,0)(0,8.125){4}{\line(-1,0){1.2}}
  \put(-1.8,-0.8){\makebox(0,0)[r]{\scriptsize $-5$}}
  \put(-1.8,7.3){\makebox(0,0)[r]{\scriptsize $0$}}
  \put(-1.8,15.5){\makebox(0,0)[r]{\scriptsize $+5$}}
  \put(-1.8,23.6){\makebox(0,0)[r]{\scriptsize $+10$}}
  \put(-6.8,13.5){\makebox(0,0)[c]{\rotatebox{90}{\scriptsize difference, \%}}}
  % x ticks
  \multiput(0,0)(20,0){3}{\line(0,-1){1.2}}
  \put(0,-3.3){\makebox(0,0)[c]{\scriptsize $10^{-4}$}}
  \put(20,-3.3){\makebox(0,0)[c]{\scriptsize $10^{-3}$}}
  \put(40,-3.3){\makebox(0,0)[c]{\scriptsize $10^{-2}$}}
  \put(20,-5.7){\makebox(0,0)[c]{\scriptsize difference step, m of control travel}}
  \put(1.5,16.9){\makebox(0,0)[l]{\scriptsize acceptance band, $\pm5\,\%$}}
  % leading-edge thickness control -- step sensitive
  \linethickness{0.5pt}
  {\color{warm}\qbezier(0,7.68)(10,4.74)(20,1.79)
   \qbezier(20,1.79)(30,13.29)(40,24.79)
   \put(0,7.68){\circle*{1.1}}\put(20,1.79){\circle*{1.1}}\put(40,24.79){\circle*{1.1}}
   \put(2,23.5){\makebox(0,0)[l]{\scriptsize leading-edge thickness}}}
  % the other three controls
  {\color{accent}\qbezier(0,9.07)(10,8.63)(20,8.20)
   \qbezier(20,8.20)(30,8.60)(40,8.99)
   \qbezier(0,9.14)(10,8.66)(20,8.19)
   \qbezier(20,8.19)(30,8.23)(40,8.28)
   \qbezier(0,10.15)(10,9.25)(20,8.36)
   \qbezier(20,8.36)(30,8.42)(40,8.48)
   \put(0,9.07){\circle*{0.9}}\put(20,8.20){\circle*{0.9}}\put(40,8.99){\circle*{0.9}}
   \put(0,9.14){\circle*{0.9}}\put(20,8.19){\circle*{0.9}}\put(40,8.28){\circle*{0.9}}
   \put(0,10.15){\circle*{0.9}}\put(20,8.36){\circle*{0.9}}\put(40,8.48){\circle*{0.9}}
   \put(5,11.9){\makebox(0,0)[l]{\scriptsize the other three controls}}}
\end{picture}

\vspace{0.5pt}
{\scriptsize The comparison is made at the middle step, chosen before the run.
Three controls sit flat across the whole range; the leading-edge thickness
control is step-sensitive, which is why it is the largest entry in Table 3.}
\end{minipage}

% --------------- checked + table 4  |  caveat box + figure 3 ----------------
\vspace{2pt}
{\color{rulegrey}\rule{\textwidth}{0.5pt}}\\[1pt]
\begin{minipage}[t]{0.605\textwidth}
\textbf{What was checked, and what was not}

\vspace{1pt}
{\footnotesize
$\bullet$ \textbf{The drag derivatives are verified against finite differences
of the same solver, to within $4.06\,\%$ at worst} and $0.042\,\%$ at best, on
all four checked controls, against a $5\,\%$ limit fixed before the run
(Table 3). \textbf{Lift derivatives are verified the same way to within
$1.185\,\%$ at worst}, at all three angles --- twelve further comparisons.\\
$\bullet$ \textbf{The check was proved able to fail.} The same comparison was
repeated at a deliberately unsuitable step: none of the four controls agreed
there, disagreeing by $94$ to $113\,\%$. The check is measuring the derivative,
not returning agreement by construction.\\
$\bullet$ \textbf{The multipoint sum is exact.} The derivative of the averaged
objective equals the weighted sum of the three individual drag derivatives to
better than $2\times10^{-13}$ relative, on every checked control.\\
$\bullet$ \textbf{Every operating point converged}: the solver's residual fell
below its $10^{-8}$ tolerance at all three angles, in both runs.\\
$\bullet$ \textbf{Verified; no experimental comparison is available for this
section at this condition,} so the derivative check is against an independent
numerical calculation and nothing here is validated against test data.\\
$\bullet$ \textbf{No grid-refinement study was run} (one mesh, $4\,032$ cells)
and \textbf{no optimiser has run} --- the records were searched for
optimisation-iteration history and there is none.}

\vspace{3pt}
\textbf{Table 4 --- Compute: what we said upfront, and what it cost.}
One core-minute is one processor core running for one minute.

\vspace{1.5pt}
{\fontsize{6.7}{7.9}\selectfont
\begin{tabular}{@{}l r r r r@{}}
\toprule
Stage & \begin{tabular}{@{}r@{}}Stated upfront,\\ core-min\end{tabular}
      & \begin{tabular}{@{}r@{}}Used,\\ core-min\end{tabular}
      & \begin{tabular}{@{}r@{}}Wall\\ clock, s\end{tabular}
      & \begin{tabular}{@{}r@{}}Used /\\ upfront\end{tabular} \\
\midrule
Feasibility check (ran first) & $8.0$ (upper limit) & $0.583$  & $35$  & $0.07$ \\
Derivative verification       & $22.1$ (estimate)   & $14.318$ & $859$ & $\mathbf{0.65}$ \\
\midrule
Total                         & ---                 & $14.901$ & $894$ & --- \\
\bottomrule
\end{tabular}}

\vspace{1pt}
{\scriptsize Used and wall-clock figures are read from the run records. Derived
cost \$0.0127 at \$0.0513 per core-hour --- derived, not metered.}

\vspace{1.5pt}
{\footnotesize \textbf{A short check ran before we committed the verification
budget.} It solved the three angles once each --- thirty-five seconds --- to
confirm the bracket was sensible and that your lift target sat inside it, and
cost $4.1\,\%$ of the work it protected. The dollar figure is small because
this is a two-dimensional section on $4\,032$ cells; cost scales with cell
count and with the number of operating points.}

\end{minipage}\hfill
\begin{minipage}[t]{0.375\textwidth}
\colorbox{boxgrey}{\begin{minipage}[t]{0.955\textwidth}
{\footnotesize\textbf{Read this before you use the numbers}\\[1pt]
$\bullet$ Shape optimisation runs next; today's work verifies the gradients
that drive it. No optimiser step has been taken --- there is no improved shape
and no drag saving on this sheet.\\
$\bullet$ Everything here is this aerofoil section at $10\,\mathrm{m\,s^{-1}}$
and Reynolds number $6.67\times10^{5}$, between $3.139$ and $7.139$ degrees ---
nothing outside that range, and nothing about stall.\\
$\bullet$ One mesh, $4\,032$ cells, so no coefficient here has a mesh error bar.\\
$\bullet$ Four of the eight surface controls were checked; no claim is made
about the other four.\\
$\bullet$ The run was serial, on one core, and a derivative verified at one
core count is a statement about that core count.\\
$\bullet$ Figure 1 is the steepest-descent direction at the current shape; the
constraints will bend the optimiser's actual step.}
\end{minipage}}

\vspace{3pt}
\textbf{Figure 3 --- Polar.} Lift and drag against angle of attack; three
solved points each.

\vspace{1.5pt}
\setlength{\unitlength}{1mm}
\begin{picture}(43,28)(-7.5,-5.5)
  \linethickness{0.4pt}
  \put(0,0){\line(1,0){26}}
  \put(0,0){\line(0,1){24}}
  \put(26,0){\line(0,1){24}}
  \multiput(5.78,0)(5.78,0){4}{\line(0,-1){1.1}}
  \put(5.78,-3.1){\makebox(0,0)[c]{\scriptsize 4}}
  \put(11.56,-3.1){\makebox(0,0)[c]{\scriptsize 5}}
  \put(17.33,-3.1){\makebox(0,0)[c]{\scriptsize 6}}
  \put(23.11,-3.1){\makebox(0,0)[c]{\scriptsize 7}}
  \put(13,-5.5){\makebox(0,0)[c]{\scriptsize $\alpha$, deg}}
  % left axis C_L 0.30..0.70
  \multiput(0,0)(0,6){5}{\line(-1,0){1.1}}
  \put(-1.7,-0.8){\makebox(0,0)[r]{\scriptsize 0.30}}
  \put(-1.7,11.2){\makebox(0,0)[r]{\scriptsize 0.50}}
  \put(-1.7,23.2){\makebox(0,0)[r]{\scriptsize 0.70}}
  \put(-7.0,12){\makebox(0,0)[c]{\rotatebox{90}{\scriptsize $C_L$, ---}}}
  % right axis C_D 0.015..0.030
  \multiput(26,0)(0,8){4}{\line(1,0){1.1}}
  \put(27.7,-0.8){\makebox(0,0)[l]{\scriptsize 0.015}}
  \put(27.7,7.2){\makebox(0,0)[l]{\scriptsize 0.020}}
  \put(27.7,15.2){\makebox(0,0)[l]{\scriptsize 0.025}}
  \put(27.7,23.2){\makebox(0,0)[l]{\scriptsize 0.030}}
  % C_L, solid, filled circles
  \linethickness{0.5pt}
  \qbezier(0.80,0.71)(6.6,6.3)(12.36,11.93)
  \qbezier(12.36,11.93)(18.1,16.9)(23.92,21.84)
  \put(0.80,0.71){\circle*{1.1}}\put(12.36,11.93){\circle*{1.1}}
  \put(23.92,21.84){\circle*{1.1}}
  \put(9.5,14.0){\makebox(0,0)[r]{\scriptsize $C_L$}}
  % C_D, dashed, open squares
  \multiput(0.80,3.58)(1.284,0.653){9}{\rule{0.9mm}{0.4pt}}
  \multiput(12.36,9.46)(1.284,1.130){9}{\rule{0.9mm}{0.4pt}}
  \put(0.10,2.88){\framebox(1.4,1.4){}}
  \put(11.66,8.76){\framebox(1.4,1.4){}}
  \put(23.22,18.93){\framebox(1.4,1.4){}}
  \put(24.5,12.0){\makebox(0,0)[r]{\scriptsize $C_D$}}
\end{picture}

\vspace{0.5pt}
{\scriptsize Error bars are smaller than the markers (Table 2). No experimental
data exists for this section at this condition, so no reference curve is
drawn.}

\end{minipage}

% ------------------------------------------------------------- voices -------
\vspace{2pt}
{\color{rulegrey}\rule{\textwidth}{0.5pt}}

\vspace{1pt}
{\footnotesize
\rolesig{Lead Researcher.} The change worth your attention is the sweep
variable. Three Mach points at this speed would have been three copies of one
answer; angle of attack is where the section actually has to compromise, and
Table 2 shows why. One shape must serve all three angles, so the objective is
their equally weighted average, not one point with the others left to follow.

\rolesig{Lead Engineer.} Figure 1 is the actionable picture: your section's
drag objective wants a blunter leading edge above all, then a thinner
mid-chord, then a slightly fuller trailing edge. That is the local direction,
not the answer --- the optimisation that runs next follows it under your
thickness and area constraints and returns the shape and the saving. If you
want the trade made against a lift constraint instead, say so before it starts
--- the lift derivatives are already verified, so it costs nothing to add.

\rolesig{Lead Numericist.} A gradient you have not checked is a direction you
are guessing at, so we checked all sixteen --- four surface controls on drag
and the same four on lift at each of three angles --- against independently
computed finite differences. Worst disagreement $4.06\,\%$, against a
$5\,\%$ limit written down before the run and not chosen after seeing the
answer; the check was driven at a deliberately wrong step to prove it can fail,
and it did, on all four. What has no error bar is the mesh: one level, no
refinement study. Say so and we will price one.}

\vspace{1pt}
{\scriptsize All figures on this sheet are read from run records retained on
disk; coefficients are printed to six significant figures.}

\end{document}
